\documentclass[12pt,a4paper]{article}
\usepackage[utf8]{inputenc}
\usepackage[T1]{fontenc}
\usepackage[french]{babel}
\usepackage{amsmath, amsfonts, amssymb}
\usepackage{geometry}
\usepackage{xcolor}
\usepackage{booktabs}
\usepackage{array}
\usepackage{tikz}
\usetikzlibrary{shapes.geometric, arrows.meta, positioning, calc}
\usepackage{float}
\usepackage{tcolorbox}
\usepackage{hyperref}
\usepackage{enumitem}

\geometry{hmargin=2.5cm, vmargin=2.5cm}

% ── Palette de couleurs ──────────────────────────────────────────────────────
\definecolor{bleuH1}{RGB}{41, 98, 255}
\definecolor{vertH2}{RGB}{0, 150, 100}
\definecolor{gris}{RGB}{245, 245, 245}
\definecolor{orange}{RGB}{220, 120, 0}

% ── Boîtes tcolorbox ─────────────────────────────────────────────────────────
\tcbuselibrary{skins, breakable}
\newtcolorbox{boxH1}[1]{colback=gris, colframe=bleuH1, fonttitle=\bfseries, title=#1}
\newtcolorbox{boxH2}[1]{colback=gris, colframe=vertH2, fonttitle=\bfseries, title=#1}
\newtcolorbox{boxInfo}[1]{colback=blue!4, colframe=orange, fonttitle=\bfseries, title=#1}

% ── Styles TikZ ──────────────────────────────────────────────────────────────
\tikzset{
  nStart/.style ={rectangle, rounded corners, draw=black, thick, fill=green!25, text width=2.8cm, align=center, font=\bfseries\small, minimum height=1cm},
  nDecis/.style ={diamond, aspect=2.2, draw=black, thick, fill=orange!25, text width=2.4cm, align=center, font=\bfseries\small, inner sep=1pt},
  nProc/.style  ={rectangle, draw=black, thick, fill=blue!10, text width=3.8cm, align=center, font=\small, minimum height=1cm},
  nHigh/.style  ={rectangle, draw=black, thick, fill=purple!20, text width=3.8cm, align=center, font=\bfseries\small, minimum height=1cm},
  nSim/.style   ={rectangle, draw=black, thick, fill=pink!55, text width=3.8cm, align=center, font=\bfseries\small, minimum height=1cm},
  nErr/.style   ={rectangle, rounded corners, draw=black, thick, fill=red!30, text width=2.6cm, align=center, font=\bfseries\small, minimum height=1cm},
  fl/.style     ={thick, ->, >=stealth}
}

% ── Titre ─────────────────────────────────────────────────────────────────────
\title{\textbf{Documentation algorithmique}\\[4pt]
       \large Heuristiques de réarrangement SBS/RS double profondeur\\
       H1 — Greedy Constructive \quad|\quad H2 — Lookahead}
\author{Projet SBS/RS}
\date{\today}

\begin{document}
\maketitle
\tableofcontents
\newpage

% ═══════════════════════════════════════════════════════════════════════════════
\section{Contexte et modélisation du problème}
% ═══════════════════════════════════════════════════════════════════════════════

\subsection{Structure physique de l'entrepôt}

L'entrepôt est composé de \textbf{colonnes}, chacune divisée en deux côtés
(Gauche / Droite). Chaque côté forme une \textbf{demi-colonne} offrant
deux positions de stockage empilées :

\begin{itemize}[noitemsep]
  \item \textbf{AVANT} — accessible directement par la navette.
  \item \textbf{ARRIÈRE} — accessible uniquement si la position AVANT
        de la même demi-colonne est \emph{vide} (contrainte LIFO stricte).
\end{itemize}

La navette se déplace horizontalement (axe X) pour prélever et déposer
les bacs. La métrique d'optimisation est la \textbf{distance totale
parcourue sur l'axe X}.

\subsection{Invariant Back-First}

Le générateur d'instances impose et vérifie la contrainte suivante
sur les configurations \emph{initiales} et \emph{cibles} :
si une position \textbf{AVANT} est occupée, la position \textbf{ARRIÈRE} correspondante doit l'être également.

\textbf{Nuance sur les états intermédiaires :} Bien que cet invariant soit strict à l'état initial et final, les heuristiques \textbf{peuvent temporairement briser cette règle} au cours de la résolution. Par exemple, la navette peut placer un bac temporaire (buffer) dans une position AVANT dont l'ARRIÈRE est vide, notamment pour éviter d'occuper l'ARRIÈRE si celui-ci est la cible d'un autre bac. Par conséquent, il n'y a jamais besoin de mécanisme complexe "d'ouvreur" car l'algorithme dépose directement là où l'impact est le moindre.

\subsection{Objectif}

À partir d'une \emph{configuration initiale} donnée, amener chaque bac
à sa \emph{position cible} en minimisant la \textbf{distance totale
parcourue par la navette sur l'axe X}, tout en respectant à chaque
instant la contrainte LIFO et l'invariant Back-First.

% ═══════════════════════════════════════════════════════════════════════════════
\section{Infrastructure commune : le graphe de dépendances}
% ═══════════════════════════════════════════════════════════════════════════════

\subsection{Définition}

À chaque itération, les deux heuristiques construisent un \textbf{graphe orienté des dépendances} :

\begin{itemize}[noitemsep]
  \item Chaque \textbf{nœud} représente un bac présent dans l'entrepôt.
  \item Une \textbf{flèche dirigée} d'un bac X vers un bac Y signifie : \emph{« Le bac X doit quitter sa position avant que le bac Y puisse atteindre sa cible. »}
\end{itemize}

Un bac sans arc entrant (degré entrant $= 0$) est dit \textbf{libre} :
rien ne l'empêche d'être déplacé directement vers sa cible finale.

\subsection{Les trois types d'arcs}

\begin{table}[H]
\centering
\renewcommand{\arraystretch}{1.35}
\begin{tabular}{>{\bfseries}p{3.2cm} p{10.5cm}}
\toprule
Type & Condition de génération de l'arc $X \to Y$ \\
\midrule
DUR\newline(physique) &
  $X$ est en position AVANT et $Y$ en position ARRIÈRE de la
  \emph{même} demi-colonne. $X$ bloque physiquement l'accès à $Y$
  (contrainte LIFO). \\
\addlinespace
OCCUPATION\newline(logique) &
  $X$ occupe l'emplacement cible de $Y$, \emph{ou} $X$ occupe la
  position AVANT de la demi-colonne dont l'ARRIÈRE est la cible de $Y$.
  $Y$ ne peut pas s'y installer tant que $X$ n'est pas parti. \\
\addlinespace
SÉQUENCEMENT\newline(logique) &
  $Y$ cible la position AVANT d'une demi-colonne, mais $X$ — qui cible
  l'ARRIÈRE de cette même demi-colonne — n'y est pas encore. $X$ doit
  être inséré en premier pour que la configuration finale soit
  physiquement valide (invariant Back-First). \\
\bottomrule
\end{tabular}
\caption{Types d'arcs du graphe de dépendances}
\end{table}

\subsection{Détection de cycles — algorithme DFS}

\subsubsection{Pourquoi un deadlock implique un cycle}

Lorsque \emph{tous} les bacs mal placés ont un arc entrant, aucun n'est
libre : le système est en \textbf{deadlock}. La théorie des graphes
(lemme du chemin dans un graphe fini) garantit qu'en suivant les arcs
sortants à partir de n'importe quel sommet sans jamais trouver de
débouché vers un sommet libre, on finira nécessairement par revisiter
un sommet déjà emprunté — c'est le \emph{point d'entrée} du cycle. Un groupe de bacs se bloque mutuellement en boucle.

Briser ce cycle en évacuant l'un de ces bacs bloquants vers un emplacement buffer est
la seule façon de restaurer un sommet de degré entrant nul et de relancer
le happy path.

\subsubsection{Algorithme DFS à trois couleurs}

L'implémentation repose sur un DFS classique avec un marquage tricolore
des sommets :

\begin{table}[H]
\centering
\renewcommand{\arraystretch}{1.3}
\begin{tabular}{cl}
\toprule
Couleur & Signification \\
\midrule
\textbf{Blanc} & Sommet non encore visité par le DFS. \\
\textbf{Gris}  & Sommet actuellement dans la pile d'appel (chemin en cours d'exploration). \\
\textbf{Noir}  & Sommet entièrement exploré — aucun cycle ne repart de lui. \\
\bottomrule
\end{tabular}
\caption{Marquage tricolore du DFS}
\end{table}

\begin{enumerate}
  \item \textbf{Initialisation.} Tous les sommets de $G$ sont colorés en
        \emph{blanc}. Une pile de chemin $\Pi$ et un tableau d'arcs
        $\mathcal{A}$ sont créés vides.

  \item \textbf{Appel récursif} \texttt{dfs}$(u)$ :
        \begin{enumerate}[label=\alph*.]
          \item Colorier $u$ en \emph{gris} ; empiler $u$ dans $\Pi$.
          \item Pour chaque arc sortant $u \to v$ dans le graphe $G$ :
                \begin{itemize}[noitemsep]
                  \item \textbf{$v$ blanc} : $v$ non visité → empiler l'arc et
                        appel récursif \texttt{dfs}$(v)$. Si un cycle est
                        retourné, propager immédiatement.
                  \item \textbf{$v$ gris} : $v$ est déjà dans la pile courante
                        → \textbf{cycle détecté !} Extraire de $\Pi$ la
                        sous-séquence depuis $v$ jusqu'au sommet courant $u$ :
                        c'est le cycle fermé $v \to \cdots \to u \to v$.
                        Retourner cette liste d'arcs.
                  \item \textbf{$v$ noir} : déjà prouvé acyclique, ignorer.
                \end{itemize}
          \item Colorier $u$ en \emph{noir} ; dépiler $u$ de $\Pi$.
        \end{enumerate}

  \item \textbf{Itération.} L'étape 2 est lancée depuis chaque sommet
        encore \emph{blanc} jusqu'à ce qu'un cycle soit trouvé (retour
        immédiat) ou que tout $G$ soit parcouru. En situation de deadlock
        strict, un cycle est \emph{toujours} trouvé.
\end{enumerate}

\subsubsection{Exemple concret — cycle à trois bacs}

Supposons trois bacs $A$, $B$, $C$ avec les dépendances suivantes :

\begin{itemize}[noitemsep]
  \item $A \to B$ : arc DUR ($A$ est en AVANT, $B$ en ARRIÈRE de la
        même demi-colonne).
  \item $B \to C$ : arc OCCUPATION ($B$ occupe la cible de $C$).
  \item $C \to A$ : arc SÉQUENCEMENT ($A$ doit être posé en ARRIÈRE
        avant que $C$ puisse occuper l'AVANT correspondant).
\end{itemize}

Aucun des trois bacs n'a de degré entrant nul : deadlock complet.
La trace du DFS est la suivante :

\begin{table}[H]
\centering
\renewcommand{\arraystretch}{1.3}
\begin{tabular}{c >{\bfseries}c c l}
\toprule
Étape & Appel & Pile $\Pi$ & Action \\
\midrule
1 & dfs($A$) & $[A]$ & $A$ colorié gris \\
2 & dfs($B$) & $[A, B]$ & $B$ colorié gris (successeur de $A$) \\
3 & dfs($C$) & $[A, B, C]$ & $C$ colorié gris (successeur de $B$) \\
4 & — & $[A, B, C]$ & Successeur de $C$ est $A$ → $A$ est \textbf{gris} \\
5 & — & — & \textbf{Cycle détecté :} extraire $[A \to B \to C \to A]$ \\
\bottomrule
\end{tabular}
\caption{Trace DFS sur le cycle $A \to B \to C \to A$}
\end{table}

Bien que tous les bacs du cycle ne soient pas accessibles directement (par exemple, $B$ est bloqué à l'ARRIÈRE), l'algorithme va identifier l'ensemble des candidats possibles à extraire, appelés les \emph{racines physiques}. Puisque $A$ est en position AVANT, il est par définition un candidat direct. Selon les coûts ou les simulations, l'heuristique (H1 ou H2) choisira le meilleur bac à déplacer. Pour illustrer la brisure de cycle, \textbf{nous supposerons pour cet exemple que l'heuristique choisit de déplacer $A$} vers un emplacement buffer.

\begin{figure}[H]
\centering
\begin{tikzpicture}[
  bn/.style   ={circle, draw=black!70, line width=1.2pt, fill=blue!12,
                font=\bfseries\large, minimum size=1.4cm},
  bufn/.style ={rectangle, rounded corners=5pt, draw=orange!80!black,
                line width=1.2pt, fill=orange!20, font=\small\bfseries,
                minimum height=1.0cm, minimum width=2.0cm, align=center},
  ar/.style   ={->, >=stealth, line width=1.1pt},
  bk/.style   ={->, >=stealth, line width=1.1pt, dashed, draw=vertH2},
  lbl/.style  ={font=\small, fill=white, inner sep=2pt, rounded corners=2pt}
]

%% ── Nœuds (layout rectangle) ──────────────────────────────────────────────
%%    A = coin supérieur gauche   B = coin supérieur droit
%%    BUF = coin inférieur gauche C = coin inférieur droit
\node (A)   [bn]    at (0,   0)  {$A$};
\node (B)   [bn]    at (5,   0)  {$B$};
\node (C)   [bn]    at (5,  -3)  {$C$};
\node (BUF) [bufn]  at (0,  -3)  {Buffer};

%% ── Arcs du cycle ──────────────────────────────────────────────────────────
\draw[ar] (A) -- node[lbl, above]{DUR}           (B);   %% A→B (haut)
\draw[ar] (B) -- node[lbl, right]{OCCUPATION}    (C);   %% B→C (droite)
\draw[ar] (C) -- node[lbl, below]{SÉQUENCEMENT}  (A|-C) %% C→A bas puis monte
               -- (A);

%% ── Brisure : A déposé en Buffer ────────────────────────────────────────────
\draw[bk]  (A) -- node[lbl, left]{\itshape dépose $A$} (BUF);

%% ── Annotations ──────────────────────────────────────────────────────────────
\node[font=\scriptsize\itshape, text=bleuH1, align=center]
      at (0, 0.95) {AVANT — racine physique};

\node[font=\scriptsize\itshape, text=vertH2, align=center]
      at (0, -3.7) {$\Rightarrow$ cycle brisé};

\end{tikzpicture}
\caption{Cycle $A \to B \to C \to A$ et brisure par mise en buffer de la racine physique $A$}
\end{figure}

\subsubsection{Extraction de la racine physique}

Tous les bacs du cycle ne sont pas nécessairement accessibles par la
navette : seuls ceux en position \textbf{AVANT} peuvent être prélevés
directement (contrainte LIFO). L'algorithme identifie les candidats
comme suit :

\begin{enumerate}[noitemsep]
  \item Pour chaque arc $(u, v)$ du cycle extrait :
        \begin{itemize}[noitemsep]
          \item Si $u$ est en \textbf{AVANT} → $u$ est un candidat
                racine physique direct.
          \item Si $u$ est en \textbf{ARRIÈRE} → seul le bac AVANT
                de la même demi-colonne est physiquement accessible ;
                c'est lui le candidat (même s'il n'appartient pas au
                cycle).
        \end{itemize}
  \item On obtient ainsi la liste complète des \textbf{candidats accessibles} par la navette pour briser ce cycle.
\end{enumerate}

\begin{tcolorbox}[colback=blue!5, colframe=blue!60!black, fonttitle=\bfseries, title=Spécificité technique : Égalité de traitement des bloqueurs]
Il est crucial de noter que l'algorithme \textbf{ne privilégie jamais} un bac qui appartient au cycle par rapport à un bloqueur « collatéral ». 

Lorsqu'un bac du cycle est en ARRIÈRE, c'est le bac en AVANT devant lui (qui peut ne pas faire partie du cycle) qui est extrait comme candidat. Une fois la liste des candidats constituée, l'heuristique les met tous \textbf{strictement sur le même plan d'égalité}. Elle évalue le score de mise en buffer (Distance + Pénalités) pour chacun et choisira le score minimum, sans se soucier de savoir si le bac sélectionné faisait partie du cycle initial.

\textbf{Pourquoi ce choix stratégique ?} \\
Parce que, du point de vue de la navette, le résultat est identique : enlever le bloqueur collatéral a exactement le même effet physique qu'enlever un bac en AVANT du cycle. Dans les deux cas, cela dégage l'accès au bac bloqué en ARRIÈRE. L'arc DUR disparaît, le degré entrant de ce bac passe à 0, et le \emph{Happy Path} est relancé ! L'heuristique agit donc de manière pragmatique : \emph{« Peu importe qui j'enlève, tant que ça ouvre la voie. Je vais donc choisir celui pour qui le détour coûte le moins cher. »}
\end{tcolorbox}

\subsubsection{Choix du bac à déplacer dans le cycle — H1 vs H2}

Une fois les racines physiques $\mathcal{R}$ identifiées, les deux
heuristiques divergent dans leur façon de choisir \emph{lequel} déplacer
et \emph{où} le mettre :

\begin{tcolorbox}[colback=gris, colframe=bleuH1, fonttitle=\bfseries,
                  title=H1 — Sélection par score de coût statique]
Pour chaque candidat accessible, la fonction
\texttt{trouver\_buffer} cherche le meilleur emplacement vide disponible en évaluant un score.
Ce score additionne la \textbf{distance totale} parcourue par la navette (aller vers le buffer, puis retour futur vers la cible)
et un ensemble de \textbf{pénalités élevées} (pour éviter les positions en AVANT, les cibles réservées, la création de nouvelles prisons ou les positions récentes).

On sélectionne ensuite la \textbf{combinaison globale bac-buffer} dont le
score est le plus faible. En cas d'égalité de score, on préfère le candidat
mis en buffer le moins récemment (anti-oscillation).

\textbf{Limite :} H1 compare les paires uniquement sur leur coût
immédiat — il ne peut pas anticiper si ce choix débloquera
davantage de bacs au prochain tour.
\end{tcolorbox}

\begin{tcolorbox}[colback=gris, colframe=vertH2, fonttitle=\bfseries,
                  title=H2 — Sélection par simulation exhaustive]
H2 explore \textbf{toutes les combinaisons possibles} entre chaque candidat et chaque emplacement vide accessible. Pour chaque combinaison,
elle \textbf{simule virtuellement} le déplacement,
reconstruit le graphe de dépendances sur cet avenir anticipé, et compte le \textbf{nombre de bacs qui deviendraient libres}.

Le tri final classe les combinaisons en privilégiant celles qui \textbf{débloquent le plus de bacs},
départagées par les pénalités éventuelles et l'historique des déplacements.

\textbf{Avantage :} H2 choisit la cassure de cycle qui
\emph{libère le plus de bacs} au tour suivant, pas seulement la moins
coûteuse en distance.
\end{tcolorbox}

\subsubsection{Effet du déplacement en buffer sur le graphe}

Après qu'un bac est déplacé vers un emplacement buffer :

\begin{enumerate}[noitemsep]
  \item Toutes les flèches sortant de ce bac dans le graphe sont supprimées
        (il ne bloque plus personne depuis son ancienne position).
  \item L'emplacement libéré peut devenir la cible ou le
        passage d'un autre bac → d'autres blocages peuvent
        disparaître.
  \item Au moins un bac du cycle n'a plus aucune flèche entrante →
        il devient libre et le \textbf{happy path reprend} à l'itération
        suivante.
  \item Le bac temporairement mis de côté devra être repositionné vers sa cible finale
        lors d'une prochaine itération (d'où le coût de
        \textbf{double handling}).
\end{enumerate}




% ═══════════════════════════════════════════════════════════════════════════════
\section{H1 — Heuristique Greedy Constructive}
% ═══════════════════════════════════════════════════════════════════════════════

\begin{boxH1}{Principe général de H1}
H1 est une heuristique \textbf{gloutonne} : à chaque itération, elle
consulte uniquement l'état présent et applique des \textbf{règles de
priorité statiques} pour choisir le meilleur mouvement immédiat,
sans simuler l'avenir.
\end{boxH1}

\subsection{Phase 1 — Happy Path : sélection du bac éligible}

\paragraph{Étape 1 — Filtrage.}
Le graphe est construit et seuls les bacs de \emph{degré entrant nul}
(libres) qui ne sont pas encore à leur cible sont retenus comme
\textbf{candidats éligibles}.

\paragraph{Étape 2 — Attribution d'un poids stratégique.}
Chaque candidat reçoit un poids selon son profil dans l'entrepôt.
Le poids le plus bas = la priorité la plus haute.

\begin{table}[H]
\centering
\renewcommand{\arraystretch}{1.35}
\begin{tabular}{c >{\bfseries}l p{9cm}}
\toprule
Poids & Profil & Justification \\
\midrule
1 & Libérateur &
  Le bac est à l'\textbf{AVANT} et bloque un bac \emph{mal placé} à
  l'ARRIÈRE. Le déplacer déverrouille immédiatement une action utile
  pour le tour suivant. Priorité absolue. \\
2 & Vulnérable &
  Le bac est à l'\textbf{ARRIÈRE}. Il risque d'être emprisonné si un
  autre bac vient se poser devant lui. Doit être extrait urgemment. \\
3 & Inoffensif &
  Le bac est à l'\textbf{AVANT} et le bac ARRIÈRE derrière lui est
  \emph{déjà à sa cible finale}. Le déplacer n'a aucun effet de
  déblocage sur l'entrepôt. Priorité minimale. \\
\bottomrule
\end{tabular}
\caption{Poids stratégiques H1 pour le Happy Path}
\end{table}

\paragraph{Étape 3 — Départage par distance.}
Si plusieurs bacs partagent le même poids, celui dont la position
est la \textbf{plus proche de la navette} (sur l'axe X) est sélectionné.
Ce critère minimise le déplacement à vide de la navette.

\paragraph{Étape 4 — Mouvement direct.}
Le bac sélectionné est déplacé directement vers sa \textbf{cible finale}.

\subsection{Phase 2 — Deadlock : cassure de cycle}

\paragraph{Étape 1 — Identification des racines physiques.}
Pour chaque arc du cycle détecté, si le bac source est en AVANT,
il est candidat. S'il est en ARRIÈRE, c'est le bac AVANT de sa
demi-colonne qui est retenu (seul accessible physiquement par la
navette LIFO).

\paragraph{Étape 2 — Évaluation personnalisée des buffers.}
Pour \textbf{chaque bac candidat individuellement}, la fonction
\texttt{trouver\_buffer} parcourt tous les emplacements vides
accessibles et calcule un score.

Ce score additionne :
\begin{itemize}[noitemsep]
  \item La \textbf{distance navette} : le trajet pour aller de la position actuelle au buffer, puis du buffer vers la cible finale.
  \item Des \textbf{pénalités cumulées} pour décourager les mauvais choix (positionner en AVANT, voler une cible, créer une prison, ou utiliser une position récemment bloquée par la liste tabou).
\end{itemize}

\begin{tcolorbox}[colback=gris, colframe=orange, fonttitle=\bfseries, title=Mécanisme clé : buffer personnalisé par bac]
La fonction renvoie le \textbf{meilleur buffer spécifique à ce bac}.
Pour un cycle $\{b_1, b_2, b_3\}$, on calcule trois paires distinctes :
$(b_1,\, \text{buf}_1,\, s_1)$, $(b_2,\, \text{buf}_2,\, s_2)$, $(b_3,\, \text{buf}_3,\, s_3)$.
Les emplacements choisis peuvent être identiques ou différents pour chaque candidat.
\textbf{Les pénalités ne dépendent que de l'emplacement buffer}, pas du bac.
Seule la composante de distance varie d'un candidat à l'autre
pour un même emplacement buffer. Par conséquent, choisir entre deux candidats
revient essentiellement à comparer leur distance au buffer — \textbf{c'est la limite
stratégique de H1} que H2 comble en simulant l'impact réel du mouvement.
\end{tcolorbox}

\begin{table}[H]
\centering
\renewcommand{\arraystretch}{1.3}
\begin{tabular}{r l}
\toprule
Pénalité (Big-M dynamique) & Cause \\
\midrule
$+1$                         & Buffer en position AVANT (légère préférence pour l'ARRIÈRE) \\
$+100 \times |\text{cols}|$  & Buffer = cible finale d'un autre bac (risque de blocage) \\
$+500 \times |\text{cols}|$  & Buffer en AVANT d'un bac mal placé à l'ARRIÈRE (créerait une prison) \\
$+1\,000 \times |\text{cols}|$ & Buffer interdit par la liste tabou \\
\bottomrule
\end{tabular}
\caption{Hiérarchie des pénalités pour le choix du buffer — Big-M proportionnel au nombre de colonnes}
\end{table}

La hiérarchie est intentionnelle et toujours garantie quelle que soit la taille de
l'entrepôt : la distance maximale ($2\times|\text{cols}|$) est toujours inférieure
au moindre Big-M ($100\times|\text{cols}|$).

\paragraph{Étape 3 — Sélection de la meilleure paire (bac, buffer).}
On compare les triplets $(b_i,\, \text{buf}_i,\, s_i)$ de tous les candidats
et on retient celui dont le score total est le plus bas. En cas d'égalité,
on préfère le bac que l'on a le moins récemment mis en buffer.

\paragraph{Étape 4 — Mouvement en buffer.}
Le bac candidat sélectionné est déplacé vers son buffer optimal.
Ce mouvement est comptabilisé comme \textbf{Double Handling}.


\subsection{Liste Tabou}

Après chaque déplacement, la position source du bac est
\textbf{interdite} pour ce bac pendant plusieurs itérations (plafonnées à 20).

Un plafond de 20 est nécessaire pour H1 : sans anticipation,
une mémoire longue est le seul garde-fou contre les oscillations.

\subsection{Logigramme H1}

\begin{figure}[H]
\centering
\resizebox{\textwidth}{!}{%
\begin{tikzpicture}[
  nSt/.style ={rectangle, rounded corners=4pt, draw=black, thick,
               fill=green!30, font=\bfseries\small,
               minimum width=2.6cm, minimum height=0.85cm, align=center},
  nDec/.style={diamond, aspect=2.4, draw=black, thick, fill=orange!25,
               font=\bfseries\small, align=center, inner sep=2pt},
  nPr/.style ={rectangle, draw=black, thick, fill=blue!10,
               font=\small, minimum width=4.8cm, minimum height=0.9cm,
               align=center},
  nHi/.style ={rectangle, draw=black, thick, fill=purple!20,
               font=\bfseries\small, minimum width=4.8cm, minimum height=0.9cm,
               align=center},
  nEr/.style ={rectangle, rounded corners=4pt, draw=black, thick,
               fill=red!30, font=\bfseries\small,
               minimum width=2.4cm, minimum height=0.85cm, align=center},
  fl/.style  ={thick, ->, >=stealth},
  lbl/.style ={font=\small, fill=white, inner sep=1.5pt}
]

%% ── Colonne centrale ─────────────────────────────────────────────────────
\node (S)   [nSt]  at (0,   0)   {Début};
\node (T)   [nDec] at (0,  -2)   {Tous\\placés ?};
\node (OK)  [nSt]  at ( 5.5,-2)  {Succès};
\node (G)   [nPr]  at (0,  -4.2) {Construire graphe $G$};
\node (L)   [nDec] at (0,  -6.4) {Bacs\\libres ?};
\node (TAB) [nPr]  at (0, -14.4) {Mise à jour liste tabou};

%% ── Branche Happy Path (droite) ─────────────────────────────────────────
\node (W)  [nHi]  at ( 5.5, -6.4) {Attribuer poids strat.\\1 : Libérateur\\2 : Vulnérable\\3 : Inoffensif};
\node (D)  [nPr]  at ( 5.5, -8.8) {Départage :\\Poids MIN + Dist. navette MIN};
\node (MC) [nPr]  at ( 5.5,-10.8) {Déplacer vers cible finale};

%% ── Branche Deadlock (gauche) — garde puis cassure ───────────────────────
\node (VA)  [nDec] at (-5.5, -6.4)  {Vides\\accessibles ?};
\node (KO)  [nEr]  at (-5.5, -4.2)  {Échec\\(instance invalide)};
\node (C)   [nPr]  at (-5.5, -8.6)  {DFS : détecter cycle};
\node (R)   [nPr]  at (-5.5,-10.6)  {Extraire racines physiques\\(bacs en AVANT)};
\node (B)   [nHi]  at (-5.5,-12.6)  {Pour chaque candidat :\\Trouver buffer optimal\\(Distance + Pénalités Big-M)};
\node (MB)  [nPr]  at (-5.5,-14.4)  {Choisir paire (bac, buffer) score MIN\\Déplacer vers buffer};

%% ── Flèches tronc ────────────────────────────────────────────────────────
\draw[fl] (S) -- (T);
\draw[fl] (T) -- node[lbl, above]{Oui} (OK);
\draw[fl] (T) -- node[lbl, right]{Non} (G);
\draw[fl] (G) -- (L);

%% ── Happy Path ───────────────────────────────────────────────────────────
\draw[fl] (L) -- node[lbl, above]{Oui} (W);
\draw[fl] (W) -- (D);
\draw[fl] (D) -- (MC);

%% ── Deadlock : garde vides_accessibles ───────────────────────────────────
\draw[fl] (L)  -- node[lbl, above]{Non} (VA);
\draw[fl] (VA) -- node[lbl, right]{Non} (KO);
\draw[fl] (VA) -- node[lbl, right]{Oui} (C);
\draw[fl] (C)  -- (R);
\draw[fl] (R)  -- (B);
\draw[fl] (B)  -- (MB);

%% ── Convergence Tabou ────────────────────────────────────────────────────
\draw[fl] (MC.south) -- ++(0,-0.5) -| (TAB.east);
\draw[fl] (MB.south) -- ++(0,-0.1) -| (TAB.west);

%% ── Retour au début — rail extérieur gauche ──────────────────────────────
\draw[fl] (TAB.south) -- ++(0,-0.5)
          -- ++(-9,0)
          -- ++(0,16.4)
          -- (S.west);

\end{tikzpicture}}
\caption{Flux de décision — H1 (poids statiques)}
\end{figure}

% ═══════════════════════════════════════════════════════════════════════════════
\section{H2 — Heuristique Lookahead}
% ═══════════════════════════════════════════════════════════════════════════════

\begin{boxH2}{Principe général de H2}
H2 est une heuristique d'\textbf{anticipation} : avant de valider
tout mouvement, elle le \textbf{simule} dans un clone de l'entrepôt,
reconstruit le graphe sur ce futur virtuel, et mesure combien de bacs
seraient débloqués. Elle choisit systématiquement le mouvement qui
\textbf{maximise la réaction en chaîne}.
\end{boxH2}

\subsection{La simulation — cœur de H2}

Pour chaque mouvement candidat (un bac vers une destination) :

\begin{enumerate}[noitemsep]
  \item \textbf{Clonage} : copie de l'état de l'entrepôt.
  \item \textbf{Application virtuelle} : effectuer le mouvement
        uniquement dans la copie.
  \item \textbf{Reconstruction} : reconstruire le graphe de dépendances
        sur la configuration copiée.
  \item \textbf{Score de déblocage} : compter le nombre exact de bacs
        mal placés qui n'ont plus aucune flèche entrante dans ce nouveau graphe.
        Plus le score est grand, plus le mouvement est intéressant.
\end{enumerate}

\subsection{Phase 1 — Happy Path avec anticipation}

\paragraph{Étape 1 — Filtrage.}
Identique à H1 : bacs éligibles = degré entrant nul, pas encore
à leur cible.

\paragraph{Étape 2 — Simulation de chaque éligible.}
Pour chaque bac éligible, simuler son déplacement vers sa cible
finale et calculer le nombre de bacs qu'il libère.

\paragraph{Étape 3 — Tri et sélection.}
Les candidats sont triés en priorité par leur score de déblocage (du plus grand au plus petit).

Les poids stratégiques de H1 (1 à 3) sont conservés comme
\textbf{second critère de départage} uniquement, pour trancher
les égalités de score de déblocage.

\paragraph{Étape 4 — Mouvement direct.}
Le bac en tête de tri est déplacé vers sa cible finale.

\subsection{Phase 2 — Deadlock : exploration exhaustive}

Là où H1 évalue chaque candidat séparément avec son meilleur buffer,
H2 explore \textbf{toutes les combinaisons} possibles entre chaque candidat et chaque emplacement disponible.
Toutes les combinaisons sont \textbf{simulées} pour mesurer leur impact réel.

\begin{tcolorbox}[colback=gris, colframe=vertH2, fonttitle=\bfseries, title=Avantage H2 sur H1 en phase Deadlock]
Là où H1 choisit la combinaison bac-buffer la moins chère en \textbf{distance},
H2 choisit la combinaison qui \textbf{libère le plus de bacs} au prochain tour.
Sur un même emplacement buffer, les pénalités sont identiques pour tous les candidats :
c'est uniquement la simulation qui permet à H2 de discriminer stratégiquement
entre deux candidats équidistants du buffer.
\end{tcolorbox}

Les pénalités (liées aux cibles, prisons, et liste tabou) appliquées sont identiques à celles de H1. Toutefois, il existe une différence fondamentale de hiérarchie :
\begin{itemize}[noitemsep]
  \item Dans H1, les pénalités sont des \textbf{interdits stricts} (elles dominent le coût en distance).
  \item Dans H2, elles ne sont que des \textbf{critères de départage}. Le tri final s'effectue en privilégiant d'abord le nombre de bacs libérés (score de déblocage maximum), puis en minimisant les pénalités.
\end{itemize}

\textbf{Conséquence pratique :} Si un mouvement place un bac dans une position très pénalisée (par exemple, création d'une prison) mais que la simulation prouve qu'il libère 3 bacs (au lieu de 2 pour un mouvement non pénalisé), H2 choisira délibérément de subir la pénalité. La simulation prime sur les interdits théoriques, car H2 "voit" factuellement la rentabilité globale de ce sacrifice à court terme.

Cette exploration peut représenter plusieurs centaines de simulations
par itération sur les grandes instances, mais garantit le choix de
dégagement \textbf{optimal à 1 coup}.

\subsection{Liste Tabou}

La durée tabou est beaucoup plus courte (plafonnée à 5 itérations).

Un plafond bas suffit : l'anticipation empêche naturellement les
oscillations. Une durée longue briderait la capacité de H2 à
revisiter des positions dans les zones denses.

\subsection{Logigramme H2}

\begin{figure}[H]
\centering
\resizebox{\textwidth}{!}{%
\begin{tikzpicture}[
  nSt/.style ={rectangle, rounded corners=4pt, draw=black, thick,
               fill=green!30, font=\bfseries\small,
               minimum width=2.6cm, minimum height=0.85cm, align=center},
  nDec/.style={diamond, aspect=2.4, draw=black, thick, fill=orange!25,
               font=\bfseries\small, align=center, inner sep=2pt},
  nPr/.style ={rectangle, draw=black, thick, fill=blue!10,
               font=\small, minimum width=4.8cm, minimum height=0.9cm,
               align=center},
  nSim/.style={rectangle, draw=black, thick, fill=pink!55,
               font=\bfseries\small, minimum width=4.8cm, minimum height=0.9cm,
               align=center},
  nEr/.style ={rectangle, rounded corners=4pt, draw=black, thick,
               fill=red!30, font=\bfseries\small,
               minimum width=2.4cm, minimum height=0.85cm, align=center},
  fl/.style  ={thick, ->, >=stealth},
  lbl/.style ={font=\small, fill=white, inner sep=1.5pt}
]

%% ── Colonne centrale ─────────────────────────────────────────────────────
\node (S)   [nSt]  at (0,   0)   {Début};
\node (T)   [nDec] at (0,  -2)   {Tous\\placés ?};
\node (OK)  [nSt]  at ( 5.5,-2)  {Succès};
\node (G)   [nPr]  at (0,  -4.2) {Construire graphe $G$};
\node (L)   [nDec] at (0,  -6.4) {Bacs\\libres ?};
\node (TAB) [nPr]  at (0, -16.2) {Mise à jour liste tabou};

%% ── Branche Happy Path (droite) ─────────────────────────────────────────
\node (SIM) [nSim] at ( 5.5, -6.4)  {SIMULER déplacement\\vers cible finale};
\node (CNT) [nPr]  at ( 5.5, -8.6)  {Calculer Score =\\Nb bacs débloqués (futur)};
\node (MX)  [nPr]  at ( 5.5,-10.8)  {Départage et choix :\\1. Score MAX\\2. Poids (1-3) MIN\\3. Dist. navette MIN};
\node (MC)  [nPr]  at ( 5.5,-13.0)  {Déplacer vers cible finale};

%% ── Branche Deadlock (gauche) — garde puis exploration ───────────────────
\node (VA)  [nDec] at (-5.5, -6.4)  {Vides\\accessibles ?};
\node (KO)  [nEr]  at (-5.5, -4.2)  {Échec\\(instance invalide)};
\node (C)   [nPr]  at (-5.5, -8.6)  {DFS : détecter cycle};
\node (R)   [nPr]  at (-5.5,-10.6)  {Extraire racines physiques\\(bacs en AVANT)};
\node (XS)  [nSim] at (-5.5,-12.8)  {SIMULER toutes les\\combinaisons possibles\\(Racines $\times$ Buffers)};
\node (MXB) [nPr]  at (-5.5,-15.0)  {Départage et choix :\\1. Score déblocage MAX\\2. Pénalités buffer MIN};
\node (MB)  [nPr]  at (-5.5,-16.2)  {Déplacer candidat vers\\meilleur buffer};

%% ── Flèches tronc ────────────────────────────────────────────────────────
\draw[fl] (S) -- (T);
\draw[fl] (T) -- node[lbl, above]{Oui} (OK);
\draw[fl] (T) -- node[lbl, right]{Non} (G);
\draw[fl] (G) -- (L);

%% ── Happy Path ───────────────────────────────────────────────────────────
\draw[fl] (L)   -- node[lbl, above]{Oui} (SIM);
\draw[fl] (SIM) -- (CNT);
\draw[fl] (CNT) -- (MX);
\draw[fl] (MX)  -- (MC);

%% ── Deadlock : garde vides_accessibles ───────────────────────────────────
\draw[fl] (L)  -- node[lbl, above]{Non} (VA);
\draw[fl] (VA) -- node[lbl, right]{Non} (KO);
\draw[fl] (VA) -- node[lbl, right]{Oui} (C);
\draw[fl] (C)  -- (R);
\draw[fl] (R)  -- (XS);
\draw[fl] (XS) -- (MXB);
\draw[fl] (MXB)-- (MB);

%% ── Convergence Tabou ────────────────────────────────────────────────────
\draw[fl] (MC.south) -- ++(0,-0.5) -| (TAB.east);
\draw[fl] (MB.south) -- ++(0,-0.1) -| (TAB.west);

%% ── Retour au début — rail extérieur gauche ──────────────────────────────
\draw[fl] (TAB.south) -- ++(0,-0.5)
          -- ++(-9,0)
          -- ++(0,18.2)
          -- (S.west);

\end{tikzpicture}}
\caption{Flux de décision — H2 (anticipation Lookahead)}
\end{figure}

% ═══════════════════════════════════════════════════════════════════════════════
\section{Pipeline d'Analyse (Benchmark Visualization)}
% ═══════════════════════════════════════════════════════════════════════════════

\begin{boxInfo}{Automatisation de l'évaluation}
Le script \texttt{benchmark\_visualization.py} consolide les résultats bruts des algorithmes et génère \textbf{14 graphiques comparatifs} permettant d'analyser l'impact de chaque variable cachée (taille, taux de remplissage, et paramètre de maintien $p$).
\end{boxInfo}

\subsection{Extraction dynamique des variables}

Le générateur d'instances encode les propriétés de chaque configuration dans son nom de fichier :
\begin{center}
\texttt{inst\_c25\_t85\_uni\_p90\_r03.csv}
\end{center}
Le script d'analyse intercepte ces paramètres dynamiquement via des expressions régulières :
\begin{itemize}[noitemsep]
  \item \textbf{c} : Nombre de colonnes de l'entrepôt (scalabilité spatiale).
  \item \textbf{t} : Taux de remplissage (densité de stockage, de 50\,\% à 92\,\%).
  \item \textbf{p} : Probabilité de maintien de la profondeur (\textit{variable de recherche clé}). Un $p$ élevé signifie que l'instance cible nécessite moins de permutations AVANT/ARRIÈRE par rapport à l'état initial.
\end{itemize}

\subsection{Les 14 Métriques Graphiques}

La suite d'analyse se divise en cinq grands axes de performance pour guider nos conclusions de recherche :
\begin{enumerate}[noitemsep]
  \item \textbf{Performance globale} : Taux de réussite (G01), Distribution de distance (G02), Temps d'exécution (G04), Efficacité des mouvements (G03).
  \item \textbf{Analyse de densité} : Taux de succès (G05, G11), Distance (G13) et surcoût de Double Handling (G07) en fonction du taux de remplissage de l'entrepôt.
  \item \textbf{Analyse du paramètre $p$} : Taux de succès (G12) et Distance (G14) selon la complexité intrinsèque de permutation.
  \item \textbf{Scalabilité spatiale} : Impact du nombre de colonnes sur le succès (G10) et la distance parcourue (G06).
  \item \textbf{Efficacité énergétique} : Ratio distance par mouvement (G08), Front de Pareto combinant le coût calcul (temps CPU) et le coût énergétique (distance sur l'axe X) (G09).
\end{enumerate}

Ces graphiques nous permettent de diagnostiquer avec précision les faiblesses des heuristiques constructives (H1 et H2) avant d'y intégrer d'éventuels composants stochastiques (métaheuristiques).

% ═══════════════════════════════════════════════════════════════════════════════
\section{Étude des hyperparamètres}
% ═══════════════════════════════════════════════════════════════════════════════

Afin de garantir que les résultats de H1 et H2 ne sont pas biaisés par un mauvais réglage empirique, nous avons mis en place un script automatisé de \emph{Grid Search} testant systématiquement différentes combinaisons de paramètres sur un échantillon d'instances d'une grande complexité. 

\subsection{Pénalité de dépôt en position AVANT ($P_{\text{avant}}$)}
Dans la fonction d'évaluation des buffers, un bac déposé à l'AVANT subit une pénalité par défaut de $+1$ par rapport à un dépôt à l'ARRIÈRE. 
\begin{itemize}[noitemsep]
  \item \textbf{Sans pénalité ($+0$)} : La navette se base purement sur la distance géographique. Le taux de réussite \textbf{s'effondre drastiquement} (chute massive pour H1 et H2), prouvant que sans distinction de profondeur, l'algorithme bloque constamment l'accès aux positions ARRIÈRE vitales.
  \item \textbf{Pénalités fortes ($+2$ ou $+5$)} : H2 ne montre aucune amélioration, et H1 subit une légère baisse de son taux de réussite. Une pénalité trop dissuasive empêche l'algorithme d'utiliser un AVANT même lorsque cela est stratégique (par exemple pour éviter de voler la cible finale d'un autre bac).
  \item \textbf{Conclusion} : La valeur $+1$ s'avère être le minimum global parfait.
\end{itemize}

\subsection{Ordre des Poids Stratégiques}
L'ordre par défaut privilégie les profils selon l'échelle suivante : \textbf{1: Libérateur, 2: Vulnérable, 3: Inoffensif}.
Nous avons testé informatiquement d'intervertir l'importance de "Libérateur" et "Vulnérable" (i.e. Vulnérable devient la priorité absolue). 

\textbf{Résultat} : Le taux de réussite reste \textbf{rigoureusement identique}.
\textbf{Explication} : Un bac "Libérateur" est par nature situé à l'AVANT (puisqu'il emprisonne son arrière) alors qu'un "Vulnérable" est toujours situé à l'ARRIÈRE. Ces deux profils n'entrent que très rarement en confrontation directe stricte, et lorsque c'est le cas, la \emph{distance de la navette} (qui agit comme second critère de départage) tranche de façon identique la décision de l'heuristique avant que l'ordre du poids ne fasse la différence.

% ═══════════════════════════════════════════════════════════════════════════════
\section{Comparaison H1 / H2}
% ═══════════════════════════════════════════════════════════════════════════════

\begin{table}[H]
\centering
\renewcommand{\arraystretch}{1.4}
\begin{tabular}{>{\bfseries}p{3.8cm} p{4.8cm} p{4.8cm}}
\toprule
Aspect & H1 — Greedy & H2 — Lookahead \\
\midrule
Critère principal\newline (Happy Path) &
  Poids statique (1–4) \newline + distance navette &
  Score de déblocage simulé \newline + poids (départage) + distance \\
Cassure de cycle &
  Meilleur buffer par score \newline de pénalité &
  Exploration exhaustive \newline candidat $\times$ buffer + simulation \\
Durée tabou &
  $\min(20,\, |V|/2)$ &
  $\min(5,\, |V|/2)$ \\
Simulation &
  Aucune &
  Clone + reconstruction $G'$ \\
Complexité/itération &
  $O(|V|^2)$ &
  $O(|V|^3)$ \\
\midrule
Taux de succès & 69,9\,\% & \textbf{87,0\,\%} ($+17,1$ pts) \\
Distance axe X & 1365,8 & \textbf{961,7} ($-29,6$\,\%) \\
Double Handling & 98,2 & \textbf{37,9} ($-61,4$\,\%) \\
\bottomrule
\end{tabular}
\caption{Synthèse algorithmique et résultats du dernier lot (376 instances)}
\end{table}

H2 investit du temps CPU (simulations par clonage) pour économiser
de la distance mécanique navette. Dans un contexte d'entrepôt
automatisé, chaque déplacement physique a un coût énergétique et
un impact sur la durée de vie du matériel : ce compromis, allié à la
chute drastique du double handling ($-61,4$\,\%), est systématiquement rentable.

\end{document}
